\documentclass{article}
\usepackage{enumitem}
\usepackage{geometry}
\usepackage{xcolor}
\geometry{margin=1in}
\title{\textcolor{blue}{White Box Testing Report: MoneyPoly}}
\author{2024115019}
\date{\today}
\begin{document}
\maketitle


\section*{\textcolor{blue}{1.2 Code Quality Analysis}}

\begin{itemize}
  \item The codebase was analyzed using \texttt{pylint}.
  \item Each warning or suggestion was addressed iteratively, with a commit for each change.
  \item The following table summarizes the issues found and fixes applied:
\end{itemize}

\begin{tabular}{|c|p{7cm}|p{7cm}|}
\hline
Iteration & Issue Found & Fix Applied (Pylint Score) \\
\hline
0 & Initial codebase: many missing docstrings, style, and unused code issues & Baseline score: 8.17/10 \\
1 & Missing module docstrings in all modules & Added module docstrings to every file (score: 8.45/10) \\
2 & Line too long in board.py, cards.py & Broke long lines to meet 100 character limit (score: 8.62/10) \\
3 & Import errors (absolute vs relative) in board.py, game.py, bank.py, player.py & Changed to relative imports to fix import errors (score: 8.79/10) \\
4 & Unused code and variables in various modules & Removed unused imports and variables (score: 9.10/10) \\
5 & Missing class/function docstrings & Added class and function docstrings throughout codebase (score: 9.32/10) \\
6 & Style issues: inconsistent indentation, spacing & Fixed indentation and spacing for PEP8 compliance (score: 9.50/10) \\
7 & Logic: property mortgage check used '== True' & Simplified to 'if prop.is\_mortgaged' in board.py (score: 9.60/10) \\
8 & Import order not following PEP8 & Reordered imports in game.py and other modules (score: 9.70/10) \\
9 & Final newline missing in game.py, player.py & Added final newline to all files (score: 9.74/10) \\
10 & Exception type not specified & Used specific exception types in main.py (score: 9.80/10) \\
11 & Minor: docstring formatting, typos & Fixed typos and improved docstring formatting (score: 9.83/10) \\
12 & Unused function parameters & Removed unused parameters in some functions (score: 9.85/10) \\
13 & Consistent use of f-strings & Standardized string formatting to f-strings (score: 9.88/10) \\
14 & Minor: improved comments and clarity & Enhanced comments for maintainability (score: 9.90/10) \\
15 & Final review: ensured all pylint warnings (except design) are fixed & Codebase clean, final pylint score: 9.91/10 \\
\hline
\end{tabular}

\vspace{1em}
\noindent \textbf{Final pylint score: 9.91/10}

\vspace{1em}
\noindent \textbf{Note:} The remaining design warnings (e.g., too many attributes/arguments/branches in some classes, relative import order) were not refactored. Addressing these would require major restructuring of the codebase. The code is clean, well-documented, and highly maintainable.


\section*{\textcolor{blue}{1.3 White Box Test Cases}}

The following 25 white-box test cases were designed to cover all branches, key variable states, and relevant edge cases in the MoneyPoly codebase. Each test is explained in simple words: why it is needed, which branch or variable state it targets, the edge case it exercises, and any errors or logical issues it found.

\begin{enumerate}

%% ── SECTION 1: buy_property ──────────────────────────────────────────────────

  \item {\color{red}\textbf{Player cannot buy property if balance is insufficient}}
    \begin{itemize}
      \item \textit{Why needed:} Ensures the buy logic correctly blocks a player from purchasing a property they cannot afford.
      \item \textit{Branch covered:} \texttt{if player.balance < prop.price: return False}
      \item \textit{Key variable state:} \texttt{player.balance = 0}, \texttt{prop.price = 100}
      \item \textit{Edge case:} Zero-balance player attempts a purchase.
      \item \textit{Errors/issues found:} None after fix. (Related to Error 3.)
      \item \textit{Result:} \textbf{PASS}
    \end{itemize}

  \item {\color{red}\textbf{Player can buy property when balance exactly equals price}}
    \begin{itemize}
      \item \textit{Why needed:} Verifies that a player with exactly enough money is allowed to buy. This catches an off-by-one operator error.
      \item \textit{Branch covered:} \texttt{if player.balance >= prop.price: proceed}
      \item \textit{Key variable state:} \texttt{player.balance = 100}, \texttt{prop.price = 100}
      \item \textit{Edge case:} Exact boundary where balance equals price.
      \item \textit{Errors/issues found:} \textbf{Error 3} — original code used \texttt{<=} instead of \texttt{<}, which incorrectly blocked valid purchases when balance was exactly equal to price. Fixed in commit ``Error 3: Fixed player balance bounds and ownership check in buy\_property''.
      \item \textit{Result:} \textbf{PASS} after fix.
    \end{itemize}

  \item {\color{red}\textbf{Player cannot buy property already owned by another player}}
    \begin{itemize}
      \item \textit{Why needed:} Ensures a player cannot steal or overwrite ownership of a property already purchased.
      \item \textit{Branch covered:} \texttt{if prop.owner is not None: return False}
      \item \textit{Key variable state:} \texttt{prop.owner = other\_player}
      \item \textit{Edge case:} Attempting to buy an already-owned property.
      \item \textit{Errors/issues found:} \textbf{Error 3} — the original code was missing the ownership check entirely, allowing hostile takeover of properties. Fixed in same commit as above.
      \item \textit{Result:} \textbf{PASS} after fix.
    \end{itemize}

  \item {\color{red}\textbf{Successful purchase deducts balance and sets owner}}
    \begin{itemize}
      \item \textit{Why needed:} Validates the happy path: balance is updated, ownership is set, and property appears in player's list.
      \item \textit{Branch covered:} Full success path in \texttt{buy\_property}.
      \item \textit{Key variable state:} \texttt{player.balance = 500}, \texttt{prop.price = 100}; after buy: \texttt{balance = 400}, \texttt{prop.owner = player}.
      \item \textit{Edge case:} Normal purchase scenario.
      \item \textit{Errors/issues found:} None.
      \item \textit{Result:} \textbf{PASS}
    \end{itemize}

%% ── SECTION 2: pay_rent ──────────────────────────────────────────────────────

  \item {\color{red}\textbf{No rent collected on a mortgaged property}}
    \begin{itemize}
      \item \textit{Why needed:} Confirms that mortgaged properties are exempt from rent, protecting players under the mortgage rule.
      \item \textit{Branch covered:} \texttt{if prop.is\_mortgaged: return} early exit.
      \item \textit{Key variable state:} \texttt{prop.is\_mortgaged = True}
      \item \textit{Edge case:} Landing on a mortgaged property.
      \item \textit{Errors/issues found:} None.
      \item \textit{Result:} \textbf{PASS}
    \end{itemize}

  \item {\color{red}\textbf{Rent is transferred to the property owner}}
    \begin{itemize}
      \item \textit{Why needed:} Core financial mechanic: rent must leave the payer and arrive at the owner. If either side is wrong, the game economy breaks.
      \item \textit{Branch covered:} Main rent payment path: \texttt{payer.balance -= rent; owner.balance += rent}
      \item \textit{Key variable state:} Both players start at \$1500; after rent of \$10: payer = \$1490, owner = \$1510.
      \item \textit{Edge case:} Standard owned property with no mortgage.
      \item \textit{Errors/issues found:} \textbf{Error 4} — rent was deducted from the payer but never credited to the owner, causing a permanent money sink. Fixed in commit ``Error 4: Fixed pay\_rent to correctly transfer money to the property owner''.
      \item \textit{Result:} \textbf{PASS} after fix.
    \end{itemize}

%% ── SECTION 3: Player movement and Go salary ─────────────────────────────────

  \item {\color{red}\textbf{Landing exactly on Go awards the \$200 salary}}
    \begin{itemize}
      \item \textit{Why needed:} Ensures the Go position triggers the salary regardless of whether the player wraps around or steps directly onto it.
      \item \textit{Branch covered:} \texttt{if new\_position == 0: award salary}
      \item \textit{Key variable state:} \texttt{position = 39}, steps = 1 → new position = 0.
      \item \textit{Edge case:} Exact landing on index 0.
      \item \textit{Errors/issues found:} None (this path worked; the related bug was passing Go without landing).
      \item \textit{Result:} \textbf{PASS}
    \end{itemize}

  \item {\color{red}\textbf{Passing Go (not landing) also awards salary}}
    \begin{itemize}
      \item \textit{Why needed:} Players must receive \$200 every time they cross Go, not just when they land exactly on it. This is a fundamental Monopoly rule.
      \item \textit{Branch covered:} \texttt{if new\_pos < old\_pos: award salary} (wrap case)
      \item \textit{Key variable state:} \texttt{position = 38}, steps = 6 → lands on 4, crossing 0.
      \item \textit{Edge case:} Wrap-around without landing on 0.
      \item \textit{Errors/issues found:} \textbf{Error 7} — original code checked \texttt{position == 0} only, completely missing the wrap-around case. Players crossing Go without landing on it never received their salary, causing premature bankruptcies. Fixed in commit ``Error 7: Fixed player movement to properly award Go salary on wrapping past index 0''.
      \item \textit{Result:} \textbf{PASS} after fix.
    \end{itemize}

  \item {\color{red}\textbf{Zero steps does not award Go salary}}
    \begin{itemize}
      \item \textit{Why needed:} A move of 0 should not trigger any salary, as no board traversal occurred.
      \item \textit{Branch covered:} \texttt{steps = 0}: no position change, no wrap possible.
      \item \textit{Key variable state:} \texttt{position = 0}, steps = 0; balance unchanged.
      \item \textit{Edge case:} Zero dice value (boundary).
      \item \textit{Errors/issues found:} None.
      \item \textit{Result:} \textbf{PASS}
    \end{itemize}

  \item {\color{red}\textbf{Negative steps wrap position correctly via modulo}}
    \begin{itemize}
      \item \textit{Why needed:} Some card effects move players backwards. The position must wrap around using modulo without going negative.
      \item \textit{Branch covered:} \texttt{new\_pos = (old\_pos + steps) \% BOARD\_SIZE}
      \item \textit{Key variable state:} \texttt{position = 5}, steps = -3 → new position = 2.
      \item \textit{Edge case:} Negative steps from card effects.
      \item \textit{Errors/issues found:} None — modulo handles this correctly.
      \item \textit{Result:} \textbf{PASS}
    \end{itemize}

%% ── SECTION 4: PropertyGroup colour-group bonus rent ─────────────────────────

  \item {\color{red}\textbf{Rent is NOT doubled with only partial group ownership}}
    \begin{itemize}
      \item \textit{Why needed:} Rent doubling only applies when one player owns ALL properties in a colour group. Partial ownership must stay at base rent.
      \item \textit{Branch covered:} \texttt{all\_owned\_by(player)} returns False → base rent returned.
      \item \textit{Key variable state:} Two properties in Brown group; each owned by a different player.
      \item \textit{Edge case:} Split ownership across two players.
      \item \textit{Errors/issues found:} \textbf{Error 6} — original code used \texttt{any()} instead of \texttt{all()}, so a player owning even one property in a group would trigger doubled rent. Fixed in commit ``Error 6: Fixed unmortgage flag clearing logic and all\_owned\_by grouping evaluation''.
      \item \textit{Result:} \textbf{PASS} after fix.
    \end{itemize}

  \item {\color{red}\textbf{Rent IS doubled when one player owns the full colour group}}
    \begin{itemize}
      \item \textit{Why needed:} Validates that complete group ownership correctly triggers the 2× rent multiplier.
      \item \textit{Branch covered:} \texttt{all\_owned\_by(player)} returns True → rent * 2.
      \item \textit{Key variable state:} Both Brown properties owned by Player A.
      \item \textit{Edge case:} Full monopoly on a group.
      \item \textit{Errors/issues found:} None after fix of Error 6.
      \item \textit{Result:} \textbf{PASS}
    \end{itemize}

%% ── SECTION 5: find_winner ───────────────────────────────────────────────────

  \item {\color{red}\textbf{find\_winner returns the player with the highest balance}}
    \begin{itemize}
      \item \textit{Why needed:} At game end, the richest player must win. If the wrong comparison is used, the poorest player wins.
      \item \textit{Branch covered:} \texttt{max(players, key=lambda p: p.net\_worth())}
      \item \textit{Key variable state:} Alice = \$2000, Bob = \$1500 → Alice wins.
      \item \textit{Edge case:} Standard two-player end-game.
      \item \textit{Errors/issues found:} \textbf{Error 5} — original code used \texttt{min()} instead of \texttt{max()}, causing the poorest player to be declared winner. Fixed in commit ``Error 5: Fixed find\_winner to return the wealthiest player rather than the poorest''.
      \item \textit{Result:} \textbf{PASS} after fix.
    \end{itemize}

  \item {\color{red}\textbf{find\_winner with a single remaining player}}
    \begin{itemize}
      \item \textit{Why needed:} When all but one player goes bankrupt, the survivor must be correctly identified.
      \item \textit{Branch covered:} Single-element list passed to max/min.
      \item \textit{Key variable state:} \texttt{game.players = [Solo]}
      \item \textit{Edge case:} Only one player left in the game.
      \item \textit{Errors/issues found:} None.
      \item \textit{Result:} \textbf{PASS}
    \end{itemize}

  \item {\color{red}\textbf{find\_winner with no players returns None}}
    \begin{itemize}
      \item \textit{Why needed:} Protects against a crash if \texttt{find\_winner} is called on an empty player list.
      \item \textit{Branch covered:} \texttt{if not players: return None}
      \item \textit{Key variable state:} \texttt{game.players = []}
      \item \textit{Edge case:} Empty player list (large input / unexpected action).
      \item \textit{Errors/issues found:} None.
      \item \textit{Result:} \textbf{PASS}
    \end{itemize}

%% ── SECTION 6: Mortgage / Unmortgage ─────────────────────────────────────────

  \item {\color{red}\textbf{Mortgage and unmortgage cycle updates flags and balances correctly}}
    \begin{itemize}
      \item \textit{Why needed:} Validates that the mortgage/unmortgage state toggles are correct and that the bank credits/debits the right amounts.
      \item \textit{Branch covered:} Both \texttt{mortgage\_property} and \texttt{unmortgage\_property} success paths.
      \item \textit{Key variable state:} \texttt{prop.is\_mortgaged} toggled True then False; balance adjusted both times.
      \item \textit{Edge case:} Full round-trip mortgage cycle.
      \item \textit{Errors/issues found:} \textbf{Error 6} (partial) — unmortgage was clearing the \texttt{is\_mortgaged} flag before the transaction completed, leading to incorrect state. Fixed in commit ``Error 6''.
      \item \textit{Result:} \textbf{PASS} after fix.
    \end{itemize}

  \item {\color{red}\textbf{Cannot mortgage a property you do not own}}
    \begin{itemize}
      \item \textit{Why needed:} Prevents a player from mortgaging another player's property.
      \item \textit{Branch covered:} \texttt{if prop.owner != player: return False}
      \item \textit{Key variable state:} \texttt{prop.owner = other\_player}
      \item \textit{Edge case:} Unexpected player action — mortgaging someone else's property.
      \item \textit{Errors/issues found:} None.
      \item \textit{Result:} \textbf{PASS}
    \end{itemize}

  \item {\color{red}\textbf{Cannot unmortgage without sufficient funds}}
    \begin{itemize}
      \item \textit{Why needed:} Ensures a player cannot unmortgage if they lack the capital to pay, preventing negative balances.
      \item \textit{Branch covered:} \texttt{if player.balance < cost: return False}
      \item \textit{Key variable state:} \texttt{player.balance = 10}, unmortgage cost = 55.
      \item \textit{Edge case:} Near-zero balance player attempting unmortgage.
      \item \textit{Errors/issues found:} None.
      \item \textit{Result:} \textbf{PASS}
    \end{itemize}

%% ── SECTION 7: Trade ─────────────────────────────────────────────────────────

  \item {\color{red}\textbf{Trade fails if seller does not own the property}}
    \begin{itemize}
      \item \textit{Why needed:} A seller can only trade what they own. This checks the ownership guard in the trade logic.
      \item \textit{Branch covered:} \texttt{if prop.owner != seller: return False}
      \item \textit{Key variable state:} \texttt{prop.owner = None} (unowned).
      \item \textit{Edge case:} Attempting to trade an unowned property.
      \item \textit{Errors/issues found:} None.
      \item \textit{Result:} \textbf{PASS}
    \end{itemize}

  \item {\color{red}\textbf{Trade fails if buyer cannot afford the agreed price}}
    \begin{itemize}
      \item \textit{Why needed:} Protects against trades that would push the buyer into negative balance.
      \item \textit{Branch covered:} \texttt{if buyer.balance < cash\_amount: return False}
      \item \textit{Key variable state:} \texttt{buyer.balance = 10}, trade price = 500.
      \item \textit{Edge case:} Low-balance buyer attempting an expensive trade.
      \item \textit{Errors/issues found:} None.
      \item \textit{Result:} \textbf{PASS}
    \end{itemize}

%% ── SECTION 8: Jail / Bankruptcy / Card Actions ──────────────────────────────

  \item {\color{red}\textbf{Player is eliminated after going bankrupt}}
    \begin{itemize}
      \item \textit{Why needed:} Core game rule: when a player's balance hits zero or below, they must be removed from the game.
      \item \textit{Branch covered:} \texttt{if player.balance <= 0: player.is\_eliminated = True}
      \item \textit{Key variable state:} \texttt{player.balance = 0}
      \item \textit{Edge case:} Zero balance (boundary between alive and bankrupt).
      \item \textit{Errors/issues found:} None.
      \item \textit{Result:} \textbf{PASS}
    \end{itemize}

  \item {\color{red}\textbf{Get Out of Jail Free card is applied correctly}}
    \begin{itemize}
      \item \textit{Why needed:} Checks that the card action \texttt{jail\_free} increments the player's jail card counter, allowing them to use it later.
      \item \textit{Branch covered:} \texttt{if card["action"] == "jail\_free": player.get\_out\_of\_jail\_cards += 1}
      \item \textit{Key variable state:} Player's \texttt{get\_out\_of\_jail\_cards} count before/after applying card.
      \item \textit{Edge case:} Drawing a special-action card.
      \item \textit{Errors/issues found:} None.
      \item \textit{Result:} \textbf{PASS}
    \end{itemize}

  \item {\color{red}\textbf{Landing on Income Tax tile deducts \$200}}
    \begin{itemize}
      \item \textit{Why needed:} Verifies that the income tax tile resolution correctly charges the player.
      \item \textit{Branch covered:} \texttt{if tile\_type == "income\_tax": player.balance -= INCOME\_TAX\_AMOUNT}
      \item \textit{Key variable state:} \texttt{player.position = 4} (income tax tile); balance decreases by 200.
      \item \textit{Edge case:} Landing on a fixed-cost special tile.
      \item \textit{Errors/issues found:} None.
      \item \textit{Result:} \textbf{PASS}
    \end{itemize}

  \item {\color{red}\textbf{Landing on Go To Jail tile sends player to jail}}
    \begin{itemize}
      \item \textit{Why needed:} Ensures the Go To Jail tile correctly transitions the player into the jail state.
      \item \textit{Branch covered:} \texttt{if tile\_type == "go\_to\_jail": player.in\_jail = True}
      \item \textit{Key variable state:} \texttt{player.position = 30}; after resolution \texttt{player.in\_jail = True}.
      \item \textit{Edge case:} Unexpected player movement to a special tile.
      \item \textit{Errors/issues found:} None.
      \item \textit{Result:} \textbf{PASS}
    \end{itemize}

  \item {\color{red}\textbf{Player is jailed after rolling three consecutive doubles}}
    \begin{itemize}
      \item \textit{Why needed:} Rolling doubles 3 times in a row must send the player to jail, preventing exploitation of free movement.
      \item \textit{Branch covered:} \texttt{if dice.doubles\_streak >= 3: go\_to\_jail(player)}
      \item \textit{Key variable state:} \texttt{doubles\_streak} starts at 2, increments to 3 on roll → triggers jail.
      \item \textit{Edge case:} Large input of consecutive doubles reaching the 3-roll threshold.
      \item \textit{Errors/issues found:} None.
      \item \textit{Result:} \textbf{PASS}
    \end{itemize}

\end{enumerate}

\vspace{1em}
\noindent \textbf{Summary:} All 25 test cases passed after fixing the 5 logical errors found during analysis. The test suite covers every major decision branch in the codebase, key variable states, and boundary/edge cases including zero values, exact boundary conditions, negative inputs, empty collections, and unexpected player actions.

\end{document}
